\documentclass[11pt, oneside]{article}
\usepackage{geometry}
\geometry{letterpaper}
%geometry{landscape}
%\usepackage[parfill]{parskip}
\usepackage{graphicx}
\usepackage{amsmath} % Added for mathematical symbols
\usepackage{amssymb}

\title{Parallelize Graph Calls}
\author{Adam}
%\date{}

\begin{document}
\maketitle
If we want to parallelize the extraction step, we'll need a robust deduplication approach. This messy doc lays out a naive, worst case approach in which we have to decompose our graph into disconnected sub-graphs each of approximately equal size. We then need to check each sub-graph for duplicates and then check each pair of sub-graphs for duplicates. We then need to repeat this process until we have no duplicates. This is a worst case scenario, but it gives us a starting point for thinking about the problem.

\section{Graph and Tokens Definitions}
Call our graph, $G$ with nodes, $\mathcal{N}$ and edges, $\mathcal{E}$:
\[ G = \{ \mathcal{N}, \mathcal{E} \} \]

\[ \mathcal{N} = \{n \in \text{nodes} \} \quad \mathcal{E} = \{ e \in \text{edges} \} \]

\[ \text{size}(G) = |\mathcal{N}| \]

Given a model, $M$, our token limit is $T_M$. Given a node, we count the tokens needed to represent it to the LLM (name and description) using a token function, $\mathcal{T}_M$.
\\
\\
\(\mathcal{T}_M(\mathcal{N})\) is approximately normally distributed about its mean on $\mathbb{N}$ (for nodes)

\[ \mathcal{T}_M(\mathcal{N}) \sim \mathcal{N}(\mu_\mathcal{N}, \sigma_\mathcal{N}) \]

$\therefore$ we can include \(\frac{T_M}{\mu_\mathcal{N}}\) total nodes; call this the ``Model Graph Capacity'', $C_M$.
\\
\\
Ex: For Anthropic, the token limit is \( T_A = 200k \)
\\
\\
Assume approximately 50 tokens per node, \( \mu_\mathcal{N} = 50 \), implying \( C_M = 4000 \).
\\
\\
\section{Async Deduplication}
If we go async we need to de-dupe. Let's look at a concrete example:

\[ \text{say } \text{size}(G) = 16,000 \]

\[ G = \bigoplus_{i=1}^{4} g_i \]

To de-dupe we need to:
\begin{enumerate}
    \item check each \( g_i \) individually
    \item check each pair of components \( g_i \oplus g_j \), note that due to the fact that components are disconnected, we have:
\end{enumerate}

\[ \text{size}(g_i \oplus g_j) = \text{size}(g_i) + \text{size}(g_j) \]

We need to subdivide any components \(g_i\) that were not reduced by at least a factor of 2 in step 1. Assume the worst case in which we don't see this reduction in any of our components, \(g_i\). Therefore, we subdivide each \(g_i\) into \(g_{i1}\) and \(g_{i2}\) and perform pairwise graph component deduplication:

So we get to:
\[ |\{g_i \}| = \frac{\text{size}(G)}{C_n} \]

We reduce further by at least a factor of 2 in step 2, so we have:

\[ |\{g_i g_j \}| = 2 \frac{\text{size}(G)}{C_n} \]

Putting this all together, we count the total individual calls in step (1) above, as well as all the pairwise compares to give:

\[ |\{g_i \}| + |\{g_i g_j \}| (|g_i| - 1) \]

\[ = \frac{\text{size}(G)}{C_n} + \frac{2 \text{size}(G)}{C_n^2} \left(\frac{\text{size}(G)}{C_n} - 1\right) \]

\[ = 2 \left( \frac{\text{size}(G)}{C_n} \right)^2 \]

So for our quantities,

\[ \text{size}(G) = 16k \]
\[ C_n = 4k \]

\[ \text{\# LLM  calls } = 2 \left( \frac{16000}{4000} \right)^2 = 32 \]

\end{document}
